\documentclass[a4paper,11pt]{article}
\usepackage{fontspec} % Allows using modern fonts
\usepackage{xcolor} % For defining colors
\usepackage{graphicx} % For including images
\usepackage{titling} % For customizing the title page
\usepackage{amsmath} % For advanced mathematics
\usepackage[
    colorlinks=true,
    linkcolor=,
    urlcolor=blue,
    citecolor=,
    pdfborder={0 0 1}
]{hyperref} % For links and references
\usepackage{unicode-math} % For using Unicode mathematics
\usepackage[margin=2cm]{geometry} % For adjusting margins
\usepackage[english]{babel} % For the Spanish language (Note: The option is 'english', but the comment says Spanish. The content is in Spanish.)
\usepackage{parskip} % A more ordinary distribution of paragraphs
\usepackage[backend=biber,style=numeric]{biblatex} % Easy control of bibliography

\addbibresource{report.bib}

\setmainfont{Noto Sans} % Changes the main font to Georgia
\setmathfont{Fira Math} % Changes the math font to TeX Gyre Schola Math
\setmonofont{Fira Code} % Changes the monospaced font to Cascadia Mono (Original comment said Consolas, but command is Cascadia Mono)

\begin{document}

% ============================
% Title Page
% ============================
\begin{titlepage}
    \centering
    \vspace*{2cm}

    {\Huge\bfseries \textsc{Compa - Compra}  \\ \Large{Autonomous Shopping Cart} \par}

    \vspace{0.5cm}
    {\Large\bfseries Autonomous University of Barcelona (UAB)

    School of Computer Engineering - Robotics 2025 \par}

    \vspace{0.5cm}

    \includegraphics[width=0.5\textwidth]{Logo_uab.png} % UAB Logo (must be in the same folder as the document)

    \vspace{0.5cm}
    {\large Pol Tomé,

    Adrià Fernandez Mata,

    David Madueño Noguer,

    Tiago David Nunes Rodrigues,

    Moisés Sánchez Pin
     \par}

    \vspace{1cm}
    \includegraphics[width=0.5\textwidth]{compacompra_logo.jpg} % CompaCompra Logo (must be in the same folder as the document)

    \vfill
    {\Large \today} % This would render as the current date

\end{titlepage}

% ============================
% Table of Contents
% ============================
\tableofcontents % Translates to "Table of Contents" when rendered

% ============================
% Abstract of the document
% ============================
\begin{abstract}
Compa-Compra is an autonomous shopping assistant that guides and assists the store user to complete their shopping easily and simply. The user communicates with ComCom through an application where products are specified.
\textbf{Make your list and ComCom will guide you!}

[Aquí habría que añadir unas pinceladas de los resultados y funcionalidades]
\end{abstract}

% ============================
% Introduction of the document
% ============================
\section{Introduction}
The act of shopping is a daily activity, but it is not always efficient.
Whether due to the difficulty of finding products, congestion in the
aisles, or the need to compare options, navigating the supermarket
can become a tedious task. \textbf{Compa Compra} is an autonomous assistant
designed to improve the shopping experience for any user,
regardless of their profile, age, or specific needs.

Through a multimedia application, the customer enters their shopping
list, and the ComCom robot processes the best route within the establishment
using advanced localization and path planning techniques.
Thanks to the \textbf{D* Lite} algorithm and the use of \textbf{Bluetooth} 
beacons, \textbf{Compa Compra}
optimizes shopping time and facilitates efficient access to each product.

\section{Theoretical Framework and Contextualization}

The \textbf{CompaCompra} project is based on two key areas of robotics and 
autonomous navigation: \textbf{2D localization} using trilateration and 
route planning with \textbf{D* Lite}. These concepts allow the assistant to
 guide the user within the supermarket with precision and adaptability.

\subsection{2D Localization}
One of the main challenges of autonomous navigation is \textbf{indoor 
localization}. To determine \textbf{ComCom}'s position, 2D trilateration is
 used with four fixed Bluetooth beacons in the corners of the supermarket.

The method uses the measurement of the \textbf{Received Signal Strength 
Indicator (RSSI)} from each beacon to estimate the distance between the 
robot and the fixed references. However, the RSSI signal can be affected 
by interference, so \textbf{filtering and error correction} techniques, 
such as moving average and Kalman filter, are applied.

Once the estimated position \((x, y)\) is obtained, it is adjusted to 
the nearest supermarket grid cell, defined as a mesh of \textbf{0.5m² 
per node}. This process allows for a discrete representation of the 
navigation space and facilitates route calculation.

For more information on trilateration, consult the documentation: \\
\href{https://en.wikipedia.org/wiki/Trilateration}{Trilateration on Wikipedia}
\parencite{trilateration-wikipedia}

\subsection{Path Finding Algorithm: D* Lite}
The optimal movement of the robot is managed by \textbf{D* Lite}, an 
efficient route-finding algorithm in dynamic environments. Based on 
the user's position and the products on the shopping list, the system 
calculates the ideal trajectory on the node grid.

Each node in the supermarket can be classified as:
\begin{itemize}
\item \textbf{Free:} traversable areas without obstacles.
\item \textbf{Occupied:} areas blocked by shelves, obstacles, or other 
customers' carts.
\end{itemize}

Initially, all known obstacles are considered, but the path is 
\textbf{dynamically recalculated} if \textbf{ComCom} detects 
new impediments during its journey. This allows the robot to adapt 
in real-time to changes in the environment.

For more information on \textbf{D* Lite}, consult the documentation:
\href{https://en.wikipedia.org/wiki/D*_Lite}{D* Lite on Wikipedia} 
\parencite{dstar-lite-wikipedia} % (Original link was to D*, I've pointed to D* Lite)

\section{Description of the CompaCompra System}

\subsection{General Architecture}
\begin{itemize}
\item System block diagram (integration of hardware, sensors, beacons, and app)
\item Description of the main modules
\end{itemize}

\subsection{2D Localization}

The robot's main way of moving is through trilateration helped by four beacons placed
around the area of the supermarket's shelfs, able to locate nodes which represent the items
that are to be retrieved for purchase. This setup could be seen as in invisible grid or graph
in which the robot can travel and calculate the fastest path to each of said nodes. In the 
Methodology section we will discus in detail the software implementation of this 
invisible grid.

\begin{itemize}
\item Description of the node grid (supermarket area modeled as 0.5m² nodes)
\item Function of Bluetooth beacons and RSSI-based position estimation
\item Position correction and rounding process
\end{itemize}

\subsection{Pathfinding}
\begin{itemize}
\item Details on the D* Lite algorithm and its use for optimal route calculation
\item Consideration of obstacles and dynamic updates to the route
\end{itemize}

\subsection{User Interaction}
Ease of use is fundamental in a product like Compa-Compra, pourpously made in order
to make people's life way easier. An app will make the interaction with the robot
way more easy and comfortable than a pannel inserted into the robot, as it will allow
the one controling it to simply use their phone, something everybody has and knows how to use
on the greater demography.

Compa-Compra's control unit which will have a wireless capability, will recive the info-signals 
needed to drive the robot towards the desired destination. For this, the Raspberry Pi 3-B will sufice,
as it provides enough power and features to carry all tasks in our proyect.


\begin{itemize}
\item User interface and steps to generate the shopping list
\end{itemize}

\section{Methodology}
\begin{itemize}
\item Description of the implementation of various algorithms and techniques
\item Tools and technologies used (software, hardware, communication protocols)
\item Testing and validation procedures, including simulated or real supermarket environments
\end{itemize}

\section{Analysis of Results and Evaluation}
\begin{itemize}
\item Presentation of localization test results (accuracy, error rate, response to obstacles)
\item Comparison between planned and actual trajectory
\item Use of graphs and tables for visual representation of results
\item Evaluation of user-robot interaction and feedback
\item Graphical documentation (e.g., videos, photos or environment schemes)
\end{itemize}

\section{Discussion}
\begin{itemize}
\item Interpretation of results in relation to stated objectives
\item Identification of limitations and possible error sources (e.g., RSSI signal interference, supermarket layout changes)
\item Improvement proposals and reflections on system applicability in other environments
\end{itemize}
\section{Conclusions and Recommendations}
\begin{itemize}
\item Summary of the most relevant findings
\item Impact of the Compacompra system on the shopping experience
\item Recommendations for future research or improvements in design and functionality
\end{itemize}

\section{Annexes (Optional)}
\begin{itemize}
\item Relevant source codes
\item Additional system diagrams
\item Raw data or complete test results
\item Graphical documentation (e.g., photos or environment schemes)
\end{itemize}

\printbibliography

\end{document}